\documentclass{article}
\usepackage{amsmath,amssymb,hyperref}
\newcommand{\Q}{\mathbb Q}\newcommand{\C}{\mathbb C}\newcommand{\Z}{\mathbb Z}
\begin{document}
\title{Loop--genus extension classes, formal support cones, and analytic branch obstructions}
\date{16 November 2026: continuation literature audit}
\maketitle
This is an unregistered research handoff. The conditional deductions below use named internal results; none is an independently verified proof of either complete exceptional Mathieu case, JC(2), HC(4), or the formal linear-gap statement. External review remains deferred. The current goal is complete theorem closure, not a new finite search campaign.

\section{A finite Manin-class landing for the remaining Mathieu coupling}
The internal Tate-graph theorem proves safety under eventual purity whenever the top boundary graph group vanishes. Its remaining alternative with a vanishing local unit obstruction requires a nonsplit global genus--Tate extension. The earlier abelian Manin lead can be applied to that extension itself, rather than only to the distinguished section.

Let $G/S$ be the actual canonical semiabelian family, with $0\to T\to G\to A\to0$. Make a finite cover splitting $T$, and let $r=\dim T$. Its algebraic extension class is the homomorphism
\[
 c:X^*(T)=\Z^r\longrightarrow A^\vee,\qquad c_i=c(e_i).
\]
\textbf{Conditional deduction.} Assuming the retained rank-one fixed-part/constant-type lemma and the canonical comparison $V^\vee(-2)=H_1(G,\Q)$, the following are equivalent:
\begin{enumerate}
\item $V$ is globally semisimple as a local system.
\item Every $c_i$ has constant type up to isogeny.
\item Every classical differential-arithmetic Manin class $M_{K,A^\vee}(c_i)$ vanishes, for $K=\C(S)^{\rm alg}$.
\end{enumerate}
Here the third equivalence uses the classical Manin-kernel theorem, not complex-linear injectivity for the different differential-algebraic Manin map.

Indeed Cartier duality identifies $H_1(G)^\vee(1)$ with the realization of $[X^*(T)\to A^\vee]$. Its lattice quotient is constant. The abelian part is globally semisimple by pure-variation semisimplicity, so global semisimplicity is equivalent to splitting this lattice extension. Pullback along each basis vector reduces splitting to the retained rank-one lemma. The classical Manin kernel is torsion plus the constant trace. Finite-cover reflection follows by averaging a splitting over a normal closure. No local semisimplicity assumption substitutes for global semisimplicity.

This gives a concrete sufficient landing for the existing local-unit criterion: if all these finitely many Manin classes vanish, eventual-pure safety is equivalent to the retained intrinsic local class at zero vanishing. It does not force any Manin class to vanish from purity. Nor does it identify the canonical Picard-complex quotient with an arbitrarily selected boundary Jacobian.

Caporaso--Viviani, \emph{Torelli theorem for stable curves}, JEMS 13 (2011), \S3.0--3.1.1, DOI \href{https://doi.org/10.4171/JEMS/281}{10.4171/JEMS/281}, make $c$ explicit for an actual nodal curve $D$. An oriented graph cycle $\gamma=\sum n_e e$ determines
\[
 c_D(\gamma)=\left[\sum_e n_e(t_e-s_e)\right]\in\prod_v J(D_v^\nu).
\]
The cycle condition makes the divisor degree zero on each component. Thus node-difference Abel--Jacobi sections are candidates for the actual coupling calculation. For the research application one must carry the maps and relations of the actual derived Picard construction through this formula; knowing only the graph and component genera is insufficient.

A constant non-torsion node-difference class illustrates the distinction: its semiabelian extension need not split up to isogeny as an algebraic group, yet its constant underlying local system is semisimple. Torsion of all node differences is unnecessarily strong.

\section{The pure torus has an exact divisor-matrix test}
Suppose separately that the actual semiabelian family is a torus. After a finite splitting cover write the actual distinguished section as $(u_1,\ldots,u_r)$ with $u_i\in\mathcal O(S)^*$. On the smooth proper compactification $\overline S$, form the integer matrix
\[
 R_{p i}=\operatorname{ord}_p(u_i),\qquad p\in\overline S\setminus S.
\]
Then the section has constant type if and only if $R=0$. The kernel of $R:\Z^r\to\Z^{\overline S\setminus S}$ is precisely the character lattice of constant-valued characters on this section. Proof: a character $u^a$ has no divisor if and only if it is a nonzero constant on the proper connected curve. The kernel is saturated. This works on a splitting cover of arbitrary genus and uses all punctures, including those above infinity. Averaging reflects the constant-type condition to the original base under the retained lemma.

A scalar norm only tests selected sums of columns. For $(u,u^{-1})$ those sums vanish even when the residue matrix is nonzero. This is consistent with the retained actual norm-one failure. Baker--Rabinoff, \href{https://arxiv.org/abs/1308.3864}{arXiv:1308.3864}, provide a related tropical Jacobian bridge but explicitly allow zero/pole cancellation on retraction. Their metric skeleton of a degenerating proper curve is not automatically the generic fibre's boundary graph. No tropical vanishing criterion is promoted here.

\section{Formal coefficients do not preclude finite standard-basis methods}
Bahloul--Takayama, \href{https://arxiv.org/abs/math/0412044v2}{math/0412044v2}, Theorems 2.0.4 and 3.0.2, prove rational polyhedrality of closed local Groebner fans for power-series ideals and homogeneous ideals in the order-homogenized ring of formal differential operators. Section 3.1 explicitly proves the formal division theorem. Convergence or algebraicity of the symbol is therefore not needed for that theorem. A finite ideal containing a formal coefficient is still different from the unproved good filtration of the actual apolar-height sequence.

The Bernstein--Sato existence theorem also includes formal power-series rings; see Alvarez Montaner--Jeffries--Nunez-Betancourt, \href{https://arxiv.org/abs/2106.08830}{arXiv:2106.08830}, Definition 2.1 and Theorem 3.3. Their Theorem 5.6 is an actual extension to nonprincipal ideals. None supplies a scalar time equation for the raw source, which is already disproved internally by a safe lacunary control. None alone bounds apolar heights.

Here is an exact conditional finish line, independent of those existence theorems. Let
\[
 E=\{(m,k)\in\mathbb N^2: \lambda^kP^m\ne0\},\quad
 h_m=\max\{k:(m,k)\in E\},\quad b_m=m-h_m.
\]
Assume $E$ is a finite union of translates of finitely generated additive monoids, whose nonzero generators $(a_i,d_i)$ have $a_i>0$. If $b_m\to+\infty$, then $b_m\ge\epsilon m-C$ for some $\epsilon>0,C$.

Proof: a generator with $d_i\ge a_i$ in any nonempty translated monoid would give an infinite subsequence with bounded-above deficit, a contradiction. For the finitely many remaining generators choose $\epsilon>0$ with $a_i-d_i\ge\epsilon a_i$. In each translate $(a,d)+\sum n_i(a_i,d_i)$, the deficit is at least $\epsilon m+(a-d-\epsilon a)$. Take the minimum of these finitely many constants; isolated translates are absorbed by $C$. Thus $h_m\le(1-\epsilon)m+C$.

The missing assertion is exactly that the ACTUAL nonvanishing support has such a description, or an equally strong exact filtered-module model. Saturation, arbitrary nested annihilators, a Groebner fan of an unrelated ideal, and cancellation-free outer support bounds do not establish it. In particular this does not reinstate the fixed-ideal-power model being tested by the Gaussian thread.

\section{Analytic invariants beyond the residual Milnor identity}
The Wright thread's residual identity is local and fixes a Milnor number; it does not fix the analytic type of the residual curve or the image branch. Three primary sources extend the previous differential-value lead:
\begin{itemize}
\item Senovilla-Sanz, \href{https://doi.org/10.1007/s00025-025-02524-0}{Results Math 80 (2025), 210}: differential values force some Bernstein--Sato roots for a one-Puiseux-pair cusp, with sharper statements for multiplicity at most four.
\item Guimaraes--Hernandes, \href{https://arxiv.org/abs/2508.17572v1}{2508.17572v1}: differential values and spectral numbers force roots for semiquasihomogeneous plane branches; comparison uses Milnor-algebra and saturated-Brieskorn-lattice filtrations.
\item Alberich-Carraminana--Morella, \href{https://arxiv.org/abs/2611.09883v1}{2611.09883v1}, submitted 6 November 2026: Theorem 3.0.1 recovers the first non-topological differential value from the reduced Bernstein--Sato roots for a one-Puiseux-pair branch. Corollary 3.0.2 recovers its Zariski invariant; Theorem 4.0.1 describes generic roots within each such stratum. These are preprint claims, not independently verified imports.
\end{itemize}
The new preprint also supplies controls: Example 6.0.1 and Remark 6.0.2 show that the same reduced Bernstein--Sato roots need not determine the remaining differential values or even the Tjurina number. Example 6.0.3 shows that a fixed differential semimodule need not determine all roots. Neither direction is a complete analytic classification. Here ``genus-one branch'' means one Puiseux pair, not an elliptic normalization.

The current logarithmic-form pullback is a possible landing for an actual comparison map. Root transport under the ramified finite map is NOT proved merely by an isomorphism of logarithmic-form modules. A valid new obstruction must first compare actual logarithmic connections, Brieskorn lattices, or a suitable direct image. No root-subset rule is asserted without that map; the one-Puiseux-pair results are not promoted to the multi-pair PF115 theorem.

\section{A recent global inequality examined and not imported}
Dimca--Pokora, \href{https://arxiv.org/abs/2610.28135v1}{2610.28135v1}, 20 October 2026, obtain an extra linear singularity constraint by differentiating Langer's orbifold Bogomolov--Miyaoka--Yau inequality at an INTERIOR coefficient where equality holds. Their Theorem 1.1 assumes a Miyaoka--Kobayashi curve of even degree at least eight with only ADE singularities. It is not an unconditional identity for a residual affine curve. Our actual compactification has not been shown to satisfy those hypotheses or equality. The portable technique is to vary boundary coefficients only after proving the correct global inequality, admissible interval, and an equality value; it is not an existing exclusion of the surviving Keller boundaries.

The balanced-torus HC4 result is already a direct internal proof. Flenner--Zaidenberg/DPD classification is already present in the research archive, so generic surface-action references were not sent as a new lead. The HC4 thread's ongoing Miyanishi proof audit was left undisturbed; the previously proposed parity deduction remains dependent on the exact external import and the internal structural reduction.
\end{document}
